APR 02 - CODIGO PERCEPTRON
import numpy as np
X=np.array([[0,0,1,1],[0,1,0,1]])
T=np.array([[0,0,0,1]])
M,N=X.shape
W=np.zeros((1,M)); b=0.
eta=1.
for epoch in range(30):
    s=W@X+b
    Y=(s>=0).astype(int)
    E=T-Y
    erros=np.count_nonzero(E)
    print(epoch,erros)
    if erros==0: break
    # batch usado no material:
    W=W+eta*(E@X.T)
    b=b+eta*E.sum()
print(W,b)
# Online, por exemplo:
w=np.zeros(M); bias=0.
for epoch in range(30):
    erros=0
    for i in range(N):
        x=X[:,i]; t=T[0,i]
        y=int(w@x+bias>=0)
        e=t-y
        w+=eta*e*x
        bias+=eta*e
        erros+=int(e!=0)
    if erros==0: break
# fronteira 2D:
# x2=-(w1*x1+b)/w2 (se w2!=0)
# nunca misturar dimensoes: X=(M,N).